% Part: first-order-logic
% Chapter: syntax-and-semantics
% Section: terms-formulas

\documentclass[../../../include/open-logic-section]{subfiles}

\begin{document}

\olfileid{fol}{syn}{frm}

\olsection{Termen en \printtoken{P}{formula}}

Als een taal~$\Lang L$ van de logica van de eerste orde is gekozen,
kunnen we uit de basistekens van~$\Lang L$ uitdrukkingen bouwen. Twee
belangrijke soorten uitdrukkingen zijn \emph{termen} en
\emph{!!{formula}s}.

\begin{defn}[Termen]
\ollabel{defn:terms}
De verzameling \emph{termen}~$\Trm[L]$ van~$\Lang L$ bouwen we stap
voor stap (inductief) met de volgende regels:
\begin{enumerate}
\item Iedere !!{variable} is een term.
\item Iedere !!{constant} van~$\Lang L$ is een term.
\item Als $f$ een !!{function} met $n$ argumentplaatsen is en $t_1$,
  \dots, $t_n$ termen zijn, dan is $\Atom{f}{t_1, \ldots, t_n}$ een
  term.
\tagitem{limitClause}{
  Verder is niets een term.}{}
\end{enumerate}
Een term zonder !!{variable}s heet een \emph{gesloten term}.
\end{defn}

\begin{explain}
In onze beschrijving staan de !!{constant}s als een aparte soort teken.
We hadden ze ook kunnen behandelen als !!{function}s met nul
argumentplaatsen. Regel twee van de termdefinitie was dan niet nodig.
We moeten $\Atom{f}{t_1, \ldots, t_n}$ dan gewoon lezen als alleen het
teken $f$ als $n = 0$.
\end{explain}

\begin{defn}[Formules]
\ollabel{defn:formulas}
Ook de verzameling \emph{!!{formula}s}~$\Frm[L]$ van de taal~$\Lang L$
bouwen we stap voor stap (inductief):
\begin{enumerate}
\tagitem{prvFalse}{$\lfalse$ is een atomaire !!{formula}.}{}

\tagitem{prvTrue}{$\ltrue$ is een atomaire !!{formula}.}{}

\item Als $R$ een !!{predicate} met $n$ argumentplaatsen van~$\Lang L$
  is en $t_1$, \dots, $t_n$ termen van~$\Lang L$ zijn, dan is
  $\Atom{R}{t_1,\ldots, t_n}$ een atomaire !!{formula}.

\item Als $t_1$ en $t_2$ termen van~$\Lang L$ zijn, dan is
  $\Atom{\eq}{t_1, t_2}$ een atomaire !!{formula}.

\tagitem{prvNot}{Als $!A$ !!a{formula} is, dan is $\lnot !A$
  !!a{formula}.}{}

\tagitem{prvAnd}{Als $!A$ en $!B$ !!{formula}s zijn, dan is $(!A \land
  !B)$ !!a{formula}.}{}

\tagitem{prvOr}{Als $!A$ en $!B$ !!{formula}s zijn, dan is $(!A \lor !B)$
  !!a{formula}.}{}

\tagitem{prvIf}{Als $!A$ en $!B$ !!{formula}s zijn, dan is $(!A \lif !B)$
  !!a{formula}.}{}

\tagitem{prvIff}{Als $!A$ en $!B$ !!{formula}s zijn, dan is $(!A \liff !B)$
  !!a{formula}.}{}

\tagitem{prvAll}{Als $!A$ !!a{formula} is en $x$ !!a{variable},
  dan is $\lforall[x][!A]$ !!a{formula}.}{}

\tagitem{prvEx}{Als $!A$ !!a{formula} is en $x$ !!a{variable},
  dan is $\lexists[x][!A]$ !!a{formula}.}{}

\tagitem{limitClause}{Verder is niets !!a{formula}.}{}
\end{enumerate}
\end{defn}

\begin{explain}
De definities van termen en !!{formula}s zijn \emph{inductieve
definities}: we bouwen de verzamelingen stap voor stap op. Voor de
!!{formula}s zijn er oneindig veel mogelijke stappen. In de eerste
stap noemen we alle atomaire formules formules. Dat zijn precies de
eerste gevallen van de definitie, dus de gevallen voor
\iftag{prvTrue}{$\ltrue$, }{}%
\iftag{prvFalse}{$\lfalse$, }{}%
$\Atom{R}{t_1,\dots,t_n}$ en $\Atom{\eq}{t_1,t_2}$. Een ``atomaire
!!{formula}'' is dus iedere !!{formula} van een van deze vormen.

De andere gevallen zijn regels om uit bestaande !!{formula}s nieuwe
!!{formula}s te maken. In de tweede stap passen we die regels toe op
atomaire !!{formula}s. In de derde stap gebruiken we zowel de atomaire
formules als de formules uit de tweede stap, enzovoort. Iets is een
!!{formula} als het uiteindelijk in zo'n stap wordt gemaakt. Iets
anders is geen !!{formula}.
\end{explain}

Volgens afspraak zetten we $\eq$ tussen zijn twee argumenten en laten
we de haakjes weg. Zo is $\eq[t_1][t_2]$ een afkorting van
$\Atom{\eq}{t_1,t_2}$. Ook korten we $\lnot \Atom{\eq}{t_1,t_2}$ af
als $\eq/[t_1][t_2]$. Een formule $(!B \ast !C)$ is uit $!B$, $!C$
gemaakt met een verbindingsteken~$\ast$ met twee argumentplaatsen. Het
buitenste paar haakjes laten we vaak weg; we schrijven dan alleen
$!B \ast !C$.

\begin{intro}
Sommige logicaboeken eisen dat de !!{variable}~$x$ echt in~$!A$
voorkomt voordat we
\iftag{prvEx}{$\lexists[x][!A]$ }{}%
\iftag{notprvEx,notprvAll}{}{en }%
\iftag{prvAll}{$\lforall[x][!A]$ }{}%
rekenen tot
\iftag{notprvEx,notprvAll}{!!a{formula}}{!!{formula}s}.
Wij stellen die eis niet. Dat veroorzaakt geen problemen en maakt de
behandeling eenvoudiger.
\end{intro}

\begin{tagblock}{defNot,defOr,defAnd,defIf,defIff,defTrue,defFalse,defEx,defAll}
\begin{defn}
Formules met de gedefinieerde operatoren zijn afkortingen. We spreken
het volgende af:

\begin{tagenumerate}{defTrue,defFalse,defNot,defOr,defAnd,defIf,defIff,defEx,defAll}

\tagitem{defTrue}{$\ltrue$ is een afkorting van
  \iftag{prvFalse}{$\lnot\lfalse$}{$(!A \lor \lnot !A)$ voor een
    zekere vaste atomaire !!{formula}~$!A$}.}{}

\tagitem{defFalse}{$\lfalse$ is een afkorting van
  \iftag{prvTrue}{$\lnot\ltrue$}{$(!A \land \lnot !A)$ voor een
    zekere vaste atomaire !!{formula}~$!A$}.}{}

\tagitem{defNot}{$\lnot !A$ is een afkorting van $!A \lif \lfalse$.}{}

\tagitem{defOr}{$!A \lor !B$ is een afkorting van
  \iftag{prvAnd}{$\lnot(\lnot !A \land \lnot !B)$}{$\lnot !A \lif
    !B$}.}{}

\tagitem{defAnd}{$!A \land !B$ is een afkorting van
  \iftag{prvOr}{$\lnot(\lnot !A \lor \lnot !B)$}{$\lnot (!A \lif
    \lnot !B)$}.}{}

\tagitem{defIf}{$!A \lif !B$ is een afkorting van
  \iftag{prvOr}{$\lnot !A \lor !B)$}{$\lnot (!A \land \lnot !B)$}.}{}

\tagitem{defIff}{$!A \liff !B$ is een afkorting van $(!A \lif !B) \land (!B
  \lif !A)$.}{}

\tagitem{defAll}{$\lforall[x][!A]$ is een afkorting van $\lnot\lexists[x][\lnot !A]$.}{}

\tagitem{defEx}{$\lexists[x][!A]$ is een afkorting van $\lnot\lforall[x][\lnot !A]$.}{}
\end{tagenumerate}
\end{defn}
\end{tagblock}

In een taal voor een bepaalde toepassing zetten we !!{predicate}s en
!!{function}s met twee argumentplaatsen vaak tussen de termen waarop ze
werken. In de rekenkunde schrijven we bijvoorbeeld $t_1 < t_2$ en
$(t_1 + t_2)$; in de verzamelingenleer schrijven we $t_1 \in t_2$. De
opvolgerfunctie uit de rekenkunde schrijven we volgens afspraak zelfs
\emph{na} haar argument:~$t'$. Officieel zijn dit alleen afkortingen.
Ze staan respectievelijk voor $\Atom{\Obj A^2_0}{t_1, t_2}$,
$\Obj f^2_0(t_1, t_2)$, $\Atom{\Obj A^2_0}{t_1, t_2}$ en $f^1_0(t)$.

\begin{defn}[Syntactische identiteit]
Het symbool $\ident$ betekent dat twee tekenreeksen syntactisch precies
gelijk zijn. Dus $!A \ident !B$ geldt precies wanneer $!A$ en $!B$
even lang zijn en op iedere plaats hetzelfde teken hebben.
\end{defn}

Aan beide kanten van $\ident$ mogen tekenreeksen staan die uit kleinere
reeksen zijn samengevoegd. Zo betekent $!A \ident (!B \lor !C)$ dat
de tekenreeks~$!A$ uit precies het volgende bestaat, in deze volgorde:
een openingshaakje, de tekenreeks $!B$, het teken $\lor$, de
tekenreeks~$!C$ en een sluithaakje. Dan weten we dus ook dat het eerste
teken van~$!A$ een openingshaakje is. Verder bevat $!A$ de tekenreeks
$!B$ vanaf het tweede teken, en direct na die deelreeks staat~$\lor$,
enzovoort.

Termen en !!{formula}s worden stap voor stap uit basiselementen
opgebouwd. Daarom kunnen we eigenschappen ervan bewijzen met de volgende
inductieprincipes.

\begin{lem}[\emph{Inductieprincipe voor termen}]
    \ollabel{lem:trmind}
    Neem een taal~$\Lang L$ van de logica van de eerste orde.
    Stel dat een eigenschap~$P$ aan de volgende voorwaarden voldoet:
    %
    \begin{enumerate}
        \item zij geldt voor iedere !!{variable}~$v$;
        %
        \item zij geldt voor iedere !!{constant}~$a$ van~$\Lang L$; en
        %
        \item zij geldt voor $f(t_1,\dotsc,t_n)$ zodra zij geldt voor
        $t_1$,~\dots, $t_n$ en $f$~een !!{function} met $n$
            argumentplaatsen van~$\Lang L$ is.
    \end{enumerate}
    Hierbij zijn $t_1$,~\dots, $t_n$ termen van~$\Lang{L}$. Dan geldt
    $P$ voor iedere term in~$\Trm[L]$.
\end{lem}

\begin{prob}
    Bewijs \olref[fol][syn][frm]{lem:trmind}.
\end{prob}

\begin{lem}[\emph{Inductieprincipe voor !!{formula}s}]
    \ollabel{thm:frmind}
    Neem een taal~$\Lang L$ van de logica van de eerste orde.
    Stel dat een eigenschap~$P$ voor alle atomaire !!{formula}s geldt.
    Stel ook dat zij aan de volgende voorwaarden voldoet:
    %
    \begin{enumerate}
        \tagitem{prvNot}{zij geldt voor $\lnot !A$ zodra zij
            geldt voor~$!A$;}{}
        \tagitem{prvAnd}{zij geldt voor $(!A \land !B)$ zodra zij
            geldt voor $!A$ en~$!B$;}{}
        \tagitem{prvOr}{zij geldt voor $(!A \lor !B)$ zodra zij
            geldt voor $!A$ en~$!B$;}{}
        \tagitem{prvIf}{zij geldt voor $(!A \lif !B)$ zodra zij
            geldt voor $!A$ en~$!B$;}{}
        \tagitem{prvIff}{zij geldt voor $(!A \liff !B)$ zodra zij
            geldt voor $!A$ en~$!B$;}{}
        \tagitem{prvEx}{zij geldt voor $\lexists[x][!A]$ zodra zij
            geldt voor~$!A$;}{}
        \tagitem{prvAll}{zij geldt voor $\lforall[x][!A]$ zodra zij
            geldt voor~$!A$;}{}
  \end{enumerate}
  Hierbij zijn $!A$ en $!B$ !!{formula}s van~$\Lang{L}$. Dan geldt
  $P$ voor alle formules in~$\Frm[L]$.
\end{lem}
\end{document}
